\documentclass[11pt,a4paper]{article}
\usepackage[margin=25mm]{geometry}
\usepackage{amsmath,amssymb,amsthm,hyperref}
\setlength{\emergencystretch}{2em}
\newtheorem{lemma}{Lemma}
\newtheorem{theorem}[lemma]{Theorem}
\newcommand{\gen}[1]{\langle#1\rangle}
\newcommand{\Sym}{\operatorname{Sym}}
\title{Two endpoint anchors and an internal perfect split}
\author{Human Compression project: a component of the L019 recovery}
\date{}
\begin{document}\maketitle
An endpoint anchor is a point moved by the first specified involution
and fixed by all the others; at the other end the order is reversed.
All groups below are actual permutation groups. An eligible edge joins
orbits of the actual maximal omission of its colour. We acquire the
complete unique-crossing, minimal-transitivity, orbit-budget, saturated
forest, and split-inheritance arguments, with their exact hypotheses.
In particular, no restriction or block quotient is assumed to inherit
the intersection property.

\begin{lemma}\label{fixed}
Suppose commuting groups $L,R$ have transitive product on $A$.
If $R$ fixes a point $e\in A$, then $L$ is transitive on $A$
and $R$ fixes $A$ pointwise.
\end{lemma}
\begin{proof}
Every element of the product sends $e$ to a point of $e^L$,
since $R$ fixes $e$ and commutes with $L$. Thus $e^L=A$.
Every element of $R$ fixes every point $e^l$, again by commutation.
\end{proof}

\begin{lemma}\label{odd}
Let a connected transitive string action of rank $t$ on $d$ points
have all maximal omissions intransitive. If $d$ is odd,
$2t=d+1$, and both orbits of a unique crossing admit nontrivial
block systems for its omission, that crossing is perfect.
\end{lemma}
\begin{proof}
Write their block numbers and sizes as $m_1,k_1,m_2,k_2\ge2$.
The complete orbit-budget argument gives
\[
 t\le m_1+m_2+k_1+k_2-3,
 \qquad
 2t\le d+2-D,
 \quad D=\sum_{j=1}^2(m_j-2)(k_j-2).
\]
Here $D\ge0$ and $D\equiv d\pmod2$, since $m_jk_j$
is the size of the corresponding orbit. Therefore $D\ge1$.
As $2t=d+1$, equality holds throughout and $D=1$.
For an inclusion-minimal subfamily $M$ transitive on the two block
sets, let $c_j$ be its number of point orbits on side $j$.
The inequalities proving the displayed bound are precisely
\[
 |M|\le m_1+m_2-2,\qquad c_j\le k_j,\qquad
 t-1-|M|\le c_1+c_2-2.
\]
They must all be equalities.

These are exactly the hypotheses of the constructive part of the
acquired saturated-forest proof. To identify its interface explicitly:
the block images of $M$ are the two natural transposition paths;
the crossing blocks are endpoints, so their labels occupy the two
intervals next to the crossing. Every remaining mark fixes the crossing
points, centralizes a full block path or its far-end deletion, and
therefore has trivial image on each original block set. Each point
orbit of $\gen M$ meets each block once. The selected eligible edges
of the remaining marks form two trees on these point orbits. Any
additional edge passing across a different selected colour would
contradict that colour's ambient eligibility. Thus every remaining
mark acts on only one original side. The marks of $M$ already do so.
Connectedness of the noncommutation diagram then puts all earlier
marks on one side and all later marks on the other. This is a perfect
split. The acquired proof after the equality step uses only the
three saturated inequalities above, not the former numerical
hypothesis $2t=d+2$; hence it applies unchanged here.
\end{proof}

\begin{theorem}\label{anchors}
Let $H=\gen{s_0,\ldots,s_{t-1}}$ be a transitive string action
with connected noncommutation diagram and all maximal omissions
intransitive. Suppose it has endpoint anchors at both ends and
$2t>d$, where $d$ is its degree. Then it has a perfect split.
\end{theorem}
\begin{proof}
If $2t>d+1$, the acquired boundary theorem applies. It remains
only $2t=d+1$ with $d$ odd. The paired-cut theorem excludes
all colours having two eligible edges, so choose a unique crossing
at $i$, with omission orbits $A,B$.

An endpoint unique crossing with a singleton orbit is perfect.
An interior singleton is impossible here: the first anchor $e$
is moved by $s_0$, so lies on the nonsingleton side $B$.
The commuting groups $L=H_{<i}$ and $R=H_{>i}$ have transitive
product on $B$, and $R$ fixes $e$. Lemma~\ref{fixed} makes
$R$ fix $B$ pointwise. It fixes the singleton as well, contrary
to the presence of its nonidentity marks. Thus both sides are
nontrivial unless the cut is already perfect.

If both sides admit nontrivial block systems, Lemma~\ref{odd}
applies. Otherwise a primitive side is available. The acquired
primitive-interval argument says that its active indices form an
interval next to the cut, except for the degree-two action of the
two neighbours; the latter has the acquired bound $2t\le d$ and
is excluded. Suppose the cut is not perfect. Reverse the order
if necessary so that a shared mark lies to the right. The primitive
interval then lies to the right, so name its side $B$; all marks
of $L=H_{<i}$ fix $B$, while $R=H_{>i}$ is transitive there.
If $i=0$, the first anchor is fixed by the whole omission, giving
a singleton orbit, a case already settled. If $i>0$, the first
anchor lies in $A$, because $s_0\in L$ moves it and fixes $B$.
The commuting product $LR$ is transitive on $A$ and $R$ fixes
this anchor. Lemma~\ref{fixed} implies that $R$ fixes $A$
pointwise. There is no shared right mark, a contradiction.
The reversed orientation uses the other endpoint anchor.
Every possible unique crossing has therefore supplied a perfect split.
\end{proof}

\begin{theorem}\label{internal}
Let a string C-marking generate the natural $S_n$, with $n\ge7$,
rank $r$, and $2r>n+2$. Then it has a perfect split whose index
belongs to $\{1,\ldots,r-2\}$.
\end{theorem}
\begin{proof}
The acquired initial-fracture theorem makes every maximal omission
intransitive, and the acquired strict theorem supplies a perfect
split. The noncommutation diagram is connected: its components would
generate commuting normal factors of $S_n$, whereas every nontrivial
normal subgroup contains $A_n$ and its centralizer is trivial.
If the split is internal there is nothing to prove. Reverse if needed
to put it at index $0$.

Its actual right wing has degree $n-1$ and rank $r-1$.
The acquired wing theorem proves, without using its final strict
inequality, that every nonsymmetric wing has rank at most half
its degree. Here $r-1>(n-1)/2$, so the wing is the natural
$S_{n-1}$. It has connected noncommutation diagram, inherits the
omission hypothesis, and satisfies
$2(r-1)>(n-1)+1$. The boundary theorem gives a perfect split
at some index $j\ge1$. If $j=1$, the adjacent-inheritance lemma
applies since $2(r-1)>n-2$ and supplies an internal parent split.
If $j\ge2$, the nonadjacent-inheritance lemma transfers it to
the parent, including $j=r-1$. Thus the only case left has
perfect splits at both $0$ and $r-1$.

Let their singleton endpoints be $a,z$. They are distinct:
$a$ is moved by $\rho_0$ and fixed by $\rho_{r-1}$,
whereas $z$ is moved by $\rho_{r-1}$ and fixed by $\rho_0$.
The two actual wings have rank $r-1>(n-1)/2$, so are respectively
$\Sym(\Omega\setminus\{a\})$ and
$\Sym(\Omega\setminus\{z\})$. The actual intersection property gives
\[
 T=\gen{\rho_1,\ldots,\rho_{r-2}}
   =\Sym(\Omega\setminus\{a,z\}).
\]
This proves transitivity of the middle group, rather than assuming it
from the two endpoint cuts. It has degree $d=n-2\ge5$, rank
$t=r-2$, and connected noncommutation diagram. Inside either wing,
the other endpoint remains a perfect split because $T$ is transitive
on the complement of its singleton. Applying the wing omission lemma
twice shows that all maximal omissions of $T$ are intransitive.

Write the outer crossing edges as $\{a,b\}$ and $\{c,z\}$.
Commutation of the outer marks shows that each fixes the other's
inner endpoint; in particular $b,c\notin\{a,z\}$. Every middle
mark other than $\rho_1$ fixes $b$, and every one other than
$\rho_{r-2}$ fixes $c$, by the unique-crossing relations.
Transitivity of $T$ forces the respective first and last marks to
move them. Thus they are endpoint anchors of the actual middle
action. Since $2t=2r-4>n-2=d$, Theorem~\ref{anchors} gives a
perfect split at an index $k\in\{1,\ldots,r-2\}$.

For $2\le k\le r-3$, transfer this split through the two endpoint
cuts using nonadjacent inheritance. If $k=1$, first transfer it
through the last endpoint to the wing omitting $\rho_0$;
it is nonadjacent to that endpoint since $r\ge5$.
Then use adjacent inheritance through the first endpoint, again with
$2(r-1)>n-2$. The case $k=r-2$ is its reverse. Every case gives
an internal perfect split of the parent.
\end{proof}

\paragraph{Attribution and scope.}
The internal-split conclusion is the corresponding step of Section~8
of Cameron--Fernandes--Leemans, \emph{The number of string C-groups
of high rank}, arXiv:2212.12723v2. The complete earlier rank theorems,
actual wing bounds, fracture and saturated-forest proofs remain
attributed acquired inputs. The new endpoint-anchor argument replaces
the source's exceptional no-perfect-split diagram classification at the
odd middle boundary. The use of the actual intersection property to
identify the middle symmetric group and the changed saturation
hypotheses are explicit. This does not yet prove a trivial internal
factor, the leaf-kernel conditions, or the reversible counting
correspondence. No whole-original or new-conjecture credit is granted.
\end{document}
